% GENERATED FILE. Edit canonical lessons, then run npm run generate:books.
% canonical-source-hash: fnv1a64:dfbbc99fe514a384
% canonical-lessons: BN-C421-mal, BN-C421-gudam, BN-C421-ordar, BN-C421-paikari, BN-C421-mojut-kora

\chapter{Goods and Stock}
\label{ch:bn-goods-and-stock}
\hlchaptermodality{421}

\begin{chapteropening}
\textbf{By the end of this chapter:} \emph{I can talk about goods in a warehouse, wholesale prices and placing an order.}

Complete the last lesson of chapter 421: I can talk about goods in a warehouse, wholesale prices and placing an order.
\end{chapteropening}

\section[\texorpdfstring{\bn{মাল} (māl)}{মাল (māl)}]{\bn{মাল} (māl) — goods; merchandise}

\label{lesson:BN-C421-mal}



\begin{quote}
Before the new one: say the Bengali for a village market, then the Bengali for to be sold.
\end{quote}

\subsection*{You'll want to know: \bn{মাল}}
\textbf{\bn{মাল}} — \emph{māl} — \textquotedblleft{}goods; merchandise\textquotedblright{}.

\begin{quote}
\bn{কুলিরা} \bn{ট্রাক} \bn{থেকে} \bn{মাল} \bn{নামাল।}
\end{quote}

\begin{quote}
\emph{kulirā ṭrāk theke māl nāmālô.}
\end{quote}

\begin{quote}
The porters unloaded the goods from the lorry.
\end{quote}

\bn{মালপত্র}, luggage, is built on it.

\begin{cousinweb}
\textbf{\bn{মাল}} (\emph{māl}) is usually traced to Arabic \emph{māl}, \textquotedblleft{}property, wealth, goods\textquotedblright{}, which reached Bengali through Persian.
\end{cousinweb}

\begin{grammarlens}[title={What you've built}]
One more word for getting around.
\end{grammarlens}

\subsection*{Your turn}
\begin{itemize}
\raggedright
  \item \emph{Say it:} \emph{māl}
  \item \emph{Say it:} \emph{māl}, once more
  \item \emph{Say it:} \emph{kulirā ṭrāk theke māl nāmālô.} — The porters unloaded the goods from the lorry.
  \item \emph{Recall:} say the Bengali for a village market, then the Bengali for to be sold, then say \emph{māl} again
\end{itemize}

\begin{tcolorbox}[breakable,colback=teal!4,colframe=teal!35!black,title={Before you move on}]
What does \bn{মাল} mean? (\textquotedblleft{}Goods; merchandise\textquotedblright{}.) Say it once more.
\end{tcolorbox}
\section[\texorpdfstring{\bn{গুদাম} (gudām)}{গুদাম (gudām)}]{\bn{গুদাম} (gudām) — a warehouse; a godown}

\label{lesson:BN-C421-gudam}



\begin{quote}
Before the new one: say the Bengali for to be sold, then the Bengali for goods.
\end{quote}

\subsection*{You'll want to know: \bn{গুদাম}}
\textbf{\bn{গুদাম}} — \emph{gudām} — \textquotedblleft{}a warehouse; a godown\textquotedblright{}.

\begin{quote}
\bn{দোকানদার} \bn{গুদামে} \bn{চালের} \bn{বস্তা} \bn{রাখলেন।}
\end{quote}

\begin{quote}
\emph{dokāndār gudāme chāler bôstā rākhlen.}
\end{quote}

\begin{quote}
The shopkeeper stored the sacks of rice in the warehouse.
\end{quote}

\begin{cousinweb}
\textbf{\bn{গুদাম}} (\emph{gudām}) is usually traced to Malay \emph{gudang}, \textquotedblleft{}a storehouse\textquotedblright{}, carried round the Indian Ocean by Portuguese traders; English \emph{godown} is the same word.
\end{cousinweb}

\begin{grammarlens}[title={What you've built}]
One more word for getting around.
\end{grammarlens}

\subsection*{Your turn}
\begin{itemize}
\raggedright
  \item \emph{Say it:} \emph{gudām}
  \item \emph{Say it:} \emph{gudām}, once more
  \item \emph{Say it:} \emph{dokāndār gudāme chāler bôstā rākhlen.} — The shopkeeper stored the sacks of rice in the warehouse.
  \item \emph{Recall:} say the Bengali for to be sold, then the Bengali for goods, then say \emph{gudām} again
\end{itemize}

\begin{tcolorbox}[breakable,colback=teal!4,colframe=teal!35!black,title={Before you move on}]
What does \bn{গুদাম} mean? (\textquotedblleft{}A warehouse; a godown\textquotedblright{}.) Say it once more.
\end{tcolorbox}
\section[\texorpdfstring{\bn{অর্ডার} (ôrḍār)}{অর্ডার (ôrḍār)}]{\bn{অর্ডার} (ôrḍār) — an order (for goods or food)}

\label{lesson:BN-C421-ordar}



\begin{quote}
Before the new one: say the Bengali for goods, then the Bengali for a warehouse.
\end{quote}

\subsection*{You'll want to know: \bn{অর্ডার}}
\textbf{\bn{অর্ডার}} — \emph{ôrḍār} — \textquotedblleft{}an order (for goods or food)\textquotedblright{}.

\begin{quote}
\bn{আমরা} \bn{ফোনে} \bn{খাবারের} \bn{অর্ডার} \bn{দিলাম।}
\end{quote}

\begin{quote}
\emph{āmrā phone khābārer ôrḍār dilām.}
\end{quote}

\begin{quote}
We ordered food by phone.
\end{quote}

\bn{অর্ডার} \bn{দেওয়া}, to place an order.

\begin{cousinweb}
\textbf{\bn{অর্ডার}} (\emph{ôrḍār}) is borrowed from English \emph{order}.
\end{cousinweb}

\begin{grammarlens}[title={What you've built}]
One more word for getting around.
\end{grammarlens}

\subsection*{Your turn}
\begin{itemize}
\raggedright
  \item \emph{Say it:} \emph{ôrḍār}
  \item \emph{Say it:} \emph{ôrḍār}, once more
  \item \emph{Say it:} \emph{āmrā phone khābārer ôrḍār dilām.} — We ordered food by phone.
  \item \emph{Recall:} say the Bengali for goods, then the Bengali for a warehouse, then say \emph{ôrḍār} again
\end{itemize}

\begin{tcolorbox}[breakable,colback=teal!4,colframe=teal!35!black,title={Before you move on}]
What does \bn{অর্ডার} mean? (\textquotedblleft{}An order (for goods or food)\textquotedblright{}.) Say it once more.
\end{tcolorbox}
\section[\texorpdfstring{\bn{পাইকারি} (pāikāri)}{পাইকারি (pāikāri)}]{\bn{পাইকারি} (pāikāri) — wholesale}

\label{lesson:BN-C421-paikari}



\begin{quote}
Before the new one: say the Bengali for a warehouse, then the Bengali for an order (for goods or food).
\end{quote}

\subsection*{You'll want to know: \bn{পাইকারি}}
\textbf{\bn{পাইকারি}} — \emph{pāikāri} — \textquotedblleft{}wholesale\textquotedblright{}.

\begin{quote}
\bn{দোকানদার} \bn{বড়} \bn{বাজার} \bn{থেকে} \bn{পাইকারি} \bn{দামে} \bn{মাল} \bn{কেনেন।}
\end{quote}

\begin{quote}
\emph{dokāndār bôṛo bājār theke pāikāri dāme māl kenen.}
\end{quote}

\begin{quote}
The shopkeeper buys goods at wholesale prices from the big market.
\end{quote}

The opposite is \bn{খুচরো}, retail, the same word as small change.

\begin{cousinweb}
\textbf{\bn{পাইকারি}} (\emph{pāikāri}) is usually traced to Persian \emph{paikār}, \textquotedblleft{}a trader, a retail dealer\textquotedblright{}, with the Bengali ending -\bn{ই}.
\end{cousinweb}

\begin{grammarlens}[title={What you've built}]
One more word for getting around.
\end{grammarlens}

\subsection*{Your turn}
\begin{itemize}
\raggedright
  \item \emph{Say it:} \emph{pāikāri}
  \item \emph{Say it:} \emph{pāikāri}, once more
  \item \emph{Say it:} \emph{dokāndār bôṛo bājār theke pāikāri dāme māl kenen.} — The shopkeeper buys goods at wholesale prices from the big market.
  \item \emph{Recall:} say the Bengali for a warehouse, then the Bengali for an order (for goods or food), then say \emph{pāikāri} again
\end{itemize}

\begin{tcolorbox}[breakable,colback=teal!4,colframe=teal!35!black,title={Before you move on}]
What does \bn{পাইকারি} mean? (\textquotedblleft{}Wholesale\textquotedblright{}.) Say it once more.
\end{tcolorbox}
\section[\texorpdfstring{\bn{মজুত} \bn{করা} (môjut kôrā)}{মজুত করা (môjut kôrā)}]{\bn{মজুত} \bn{করা} (môjut kôrā) — to stock up; to store}

\label{lesson:BN-C421-mojut-kora}



\begin{quote}
Before the new one: say the Bengali for an order (for goods or food), then the Bengali for wholesale.
\end{quote}

\subsection*{You'll want to know: \bn{মজুত} \bn{করা}}
\textbf{\bn{মজুত} \bn{করা}} — \emph{môjut kôrā} — \textquotedblleft{}to stock up; to store\textquotedblright{}.

\begin{quote}
\bn{শীতের} \bn{আগে} \bn{দোকানদার} \bn{কম্বল} \bn{মজুত} \bn{করলেন।}
\end{quote}

\begin{quote}
\emph{shīter āge dokāndār kômbôl môjut kôrlen.}
\end{quote}

\begin{quote}
Before winter the shopkeeper stocked up on blankets.
\end{quote}

Past (he/she …): \bn{তিনি} \bn{চাল} \bn{মজুত} \bn{করলেন}, \emph{tini chāl môjut kôrlen}.

Also said of traders who hold back goods to push up the price.

\begin{cousinweb}
\textbf{\bn{মজুত} \bn{করা}} (\emph{môjut kôrā}) joins \bn{মজুত}, \textquotedblleft{}in stock\textquotedblright{}, usually traced to Arabic \emph{mawjūd}, \textquotedblleft{}existing, at hand\textquotedblright{}, through Persian, to \bn{করা}, to do.
\end{cousinweb}

\begin{grammarlens}[title={What you've built}]
One more word for getting around.
\end{grammarlens}

\subsection*{Your turn}
\begin{itemize}
\raggedright
  \item \emph{Say it:} \emph{môjut kôrā}
  \item \emph{Say it:} \emph{môjut kôrā}, once more
  \item \emph{Say it:} \emph{shīter āge dokāndār kômbôl môjut kôrlen.} — Before winter the shopkeeper stocked up on blankets.
  \item \emph{Recall:} say the Bengali for an order (for goods or food), then the Bengali for wholesale, then say \emph{môjut kôrā} again
\end{itemize}

\begin{tcolorbox}[breakable,colback=teal!4,colframe=teal!35!black,title={Before you move on}]
What does \bn{মজুত} \bn{করা} mean? (\textquotedblleft{}To stock up; to store\textquotedblright{}.) Say it once more.
\end{tcolorbox}
